\section{Introduction}\label{sec:intro}

Central banks, finance ministries and firms in commodity-dependent economies use inflation forecasts, and a forecast is useful for decisions only together with a statement of its uncertainty. A prediction interval with a nominal level of 90 percent claims that in about nine cases out of ten the realised value falls inside it. Whether this holds is the question of calibration, and it can be checked against realised data; it is not guaranteed by the forecasting model, because the model's error distribution is estimated from a past that may not resemble the future.

Commodity-dependent economies are a demanding setting for calibration: their price levels are exposed to exchange-rate and commodity-price movements (\citealp{gelos2017commodity}), and their inflation regimes can change within a few years, for which we cite no source; it is the motivation for the setting and not a result, and the paper does not compare commodity-dependent economies with other economies (Section~\ref{sec:scope}). A calibration made once on a past window may fail later, and a method that adapts the interval width online is a natural candidate. Conformal prediction \citep{vovk2005algorithmic} gives distribution-free finite-sample coverage under exchangeability, which time series violate. Adaptive conformal methods \citep{gibbs2021aci,gibbs2024faci,angelopoulos2023pid,bhatnagar2023saocp} were designed for that case, and a recent regime-aware variant, DRACP \citep{oancea2026dracp}, reports mean coverage closest to nominal over 48 series including euro-area inflation. Conformal intervals have been applied to inflation in Brazil, Russia, India and China \citep{sengupta2025fewnet} and in Tunisia \citep{brahem2026tunisia}, but in the parts we read we found no comparison of the empirical coverage of several adaptive conformal methods. Whether DRACP-type results transfer to commodity-dependent economies and to multi-month horizons has therefore not been tested to our knowledge; this rests on a limited search (Section~\ref{sec:related}), not a systematic review.

\paragraph{Problem, gap and what is done.} The problem is the calibration of prediction intervals for the cumulative log growth of the consumer price index at horizons of 1, 3, 6 and 12 months in economies exposed to commodity and exchange-rate shocks. The gap is an evidence gap, not a methodological one. The closest work, \citet{oancea2026dracp}, evaluates DRACP on 48 macroeconomic and financial series, mainly one step ahead, and states that its advantage is specific to one-step-ahead forecasting and depends on the point forecaster. What is missing is a benchmark of the same family of methods on a different panel, at multi-month horizons with outcomes that arrive with a delay, with decision rules fixed before the test run so that null outcomes can be reported as such. We provide that benchmark. We do not propose a new interval method and claim no novelty beyond a controlled transfer of published methods to a new setting and a design that reports power and null results explicitly.

We construct a rolling-origin benchmark in which six interval methods (split, rolling, ACI, FACI, conformal PID, DRACP) are placed on top of point forecasts at nominal level 0.90, with outcome feedback delayed by $h$ months as in real use. The panels comprise 13 economies used to design the study and 7 hold-out economies for which no forecast had been computed before the pre-specified hold-out run. The primary outcome is the calibration error, defined per country as the absolute difference between empirical coverage and 0.90 and averaged over countries; pooled coverage is reported but is not primary, because opposite deviations of different countries cancel in a pooled average. The decision rules R1 to R5 were written in the repository before the hold-out run, and outcomes are reported as the rules classify them. ``Pre-specified'' refers to a self-attested commit history, not to an external registry, and the rules were written after the exploratory results were known (Section~\ref{sec:prespec}). An extension, specified before any forecast was computed for its countries, adds a second hold-out panel of 10 economies, the zero-shot foundation model Chronos as a fourth base model and rules E1--E5.

\paragraph{Research questions.} RQ1: are adaptive conformal intervals better calibrated than a conformal calibration performed once (split) and than the native intervals of the forecasting model? RQ2: do adaptive methods improve on recomputing the conformal quantile from a rolling window (rule R2), and does DRACP improve on ACI or is it equivalent to it (rule R3)? RQ3: how much of the under-coverage at 12 months is caused by the delay with which outcomes become known (rule R4)? RQ4: do the conclusions depend on the country panel, the base forecaster, the warm-up size and the use of a zero-shot foundation model?

\paragraph{Summary of findings.} Adaptive methods were better calibrated than split conformal, with a direction consistent across samples, base forecasters and leave-one-country-out removals, but the pre-specified rule R1 was met on the first hold-out panel only at the 96-month warm-up (a sensitivity run, not promoted); at warm-up 48 the contrasts were $-3.0$~pp ($p=0.054$) at $h=3$ and $-2.8$~pp ($p=0.092$) at $h=6$, with intervals that include zero. Whether ACI improves on a rolling window depends on the panel (not shown and inconclusive on two panels, met under the pre-specified rule in part of the cells of the second hold-out panel); the reason is unknown and the panels are not averaged. DRACP against ACI is inconclusive, and all four primary contrasts of the first hold-out panel are inconclusive under the power rule R5. At 12 months pooled coverage is below nominal for every method on the design and first hold-out panels, and delayed feedback explains part of the shortfall but did not reach the R4 threshold ($+2.4$~pp against 3~pp). No efficiency advantage in width or interval score is established. Chronos did not change the calibration conclusions, and on 22 clean target months its accuracy was descriptively at least that of LightGBM and ARIMA, under an assumed pretraining cutoff of 2025-01 that the model card does not state. \ExplainerStatus{}

\paragraph{What is not claimed.} This work does not show that adaptive methods are better than rolling windows, does not show that they are not (the answer depends on the panel), and does not show that DRACP fails or does not transfer. Most pre-specified rules returned null or inconclusive outcomes, and a null outcome here is not evidence of equality: a contrast whose minimum detectable difference is undefined, or above 7~pp, is called inconclusive and never ``no effect''. The hold-out countries are untouched with respect to countries only (Section~\ref{sec:boundary}).

\paragraph{Contributions.} (i) A calibration benchmark on 30 economies in three panels with power-aware inference, whose null outcomes were specified in advance; (ii) evidence that adaptive calibration beats one-shot calibration and, post hoc, the native ARIMA intervals; (iii) a feedback-delay diagnostic; (iv) a second hold-out panel with four base models that replicates the advantage over one-shot calibration at the primary warm-up and shows that the answer on ACI against a rolling window is panel dependent, reported and not resolved; (v) a descriptive accuracy comparison of Chronos on months it cannot be assumed to have seen; and (vi) an open pipeline with hashed snapshots, tests and a committed protocol. A platform module that serves inflation intervals is an engineering artefact and not a scientific contribution.

Section~\ref{sec:related} gives related work, Section~\ref{sec:method} the formal methods, Section~\ref{sec:data} the data and design including what was fixed before which run, Section~\ref{sec:results} the results, and Sections~\ref{sec:discussion} to~\ref{sec:repro} the discussion, limitations, conclusion and reproducibility.
